\subsection{Neighbourhoods of the origin in a topological vector space over a valued division ring}

DEFINITION 3. --- Let K be a valued division ring and E a left vector space over K; we say that a subset M of E is balanced if, for all \(x \in M\) and all \(\lambda \in K\) such that \(|\lambda| \leqslant 1\) , it is true that \(\lambda x \in M\) (or in other words if \(\lambda M \subset M\) when \(|\lambda| \leqslant 1\) ).

PROPOSITION 2. --- In a topological vector space E over a valued division ring K, the closure of a balanced set M, is a balanced set.

If B is the set of \(\xi\in K\) with \(|\xi|\leqslant1\) ; then B is closed in K. But \(B\times M\) is mapped into M by the continuous mapping \((\lambda,x)\mapsto\lambda x\) ; and therefore \(B\times\overline{M}\) is mapped into \(\overline{M}\) (GT, I, §2.1, th. 1) which proves that \(\overline{M}\) is balanced.

When M is an arbitrary set in the vector space E over a valued division ring K, the set \(M_{1}\) of the \(\lambda x\) with \(x \in M\) and \(\lambda \in K\) such that \(|\lambda| \leqslant 1\) , is clearly the smallest balanced set containing M; \(M_{1}\) is called the balanced envelope of M.

PROPOSITION 3. --- Let K be a valued locally compact and non-discrete division ring and E be a Hausdorff topological vector space (resp. a topological vector space) over K. For every compact (resp. precompact) set H in E, the balanced envelope of H is compact (resp. precompact).

If B denotes the ball \(|\xi| \leqslant 1\) in K, the balanced envelope of H is \(H_{1}\) , the image of \(B \times H\) under the continuous mapping \(m: (\lambda, x) \mapsto \lambda x\) . If E is Hausdorff, if B is compact and if H is compact then so is \(B \times H\) and therefore \(H_{1}\) . If H is precompact the restriction of m to \(B \times H\) is uniformly continuous (I, p.~5, Remark 2) and as \(B \times H\) is precompact, so also is its image under m (GT, II, § 4.2, prop. 2).

Note that the balanced envelope of a closed set is not necessarily closed. For example, in \(\mathbf{R}^2\) , the balanced envelope of the hyperbola defined by the equation \(xy = 1\) is not closed.

The union of a family of balanced sets in E is balanced, which implies that for every set M of E there is a largest balanced subset N of M called the balanced core of M; also N is not empty if and only if \(0 \in M\) . To say that \(x \in N\) means that for all \(\lambda \in K\) such that \(|\lambda| \leqslant 1\) , we have \(\lambda x \in M\) , or again (if \(0 \in M\) ) that, for all \(\mu \in K\) with \(|\mu| \geqslant 1\) , we have \(x \in \mu M\) . If \(0 \in M\) , the balanced core N of M is therefore the intersection \(\bigcap_{|\mu| \geqslant 1} \mu M\) . This shows in particular that if M is closed, so also is N.

DEFINITION 4. --- Let K be a non-discrete valued division ring and E be a left vector space on K with two subsets A and B. We say that A absorbs B if there exists \(\alpha > 0\) such that \(\lambda A \supset B\) for every \(\lambda \in K\) with \(|\lambda| \geqslant \alpha\) (or equivalently if \(\mu B \subset A\) for \(\mu \neq 0\) and \(|\mu| \leqslant \alpha^{-1}\) ). A set A of E is called absorbent if it absorbs every set consisting of a single point.

Let A be a balanced set of E; for it to absorb a set B of E it is sufficient that there exists \(\lambda \neq 0\) such that \(\lambda A \supset B\) ; in fact, for \(|\mu| \geqslant |\lambda|\) , we have \(\lambda A = (\lambda \mu^{-1}) \mu A\) , and as \(\mu A\) is balanced and \(|\lambda \mu^{-1}| \leqslant 1\) , it follows that \(\lambda A \subset \mu A\) , and thus \(B \subset \mu A\) . In particular for a balanced set A of E to be absorbent, it is necessary and sufficient that for every \(x \in E\) , there exists \(\lambda \neq 0\) in K such that \(\lambda x \in A\) . Every absorbent set of E generates the vector space E. Every finite intersection of absorbent sets is an absorbent set.

PROPOSITION 4. --- In a topological vector space E on a non-discrete valued division ring K there exists a fundamental system \(\mathfrak{B}\) of closed neighbourhoods of 0 such that :

\((\mathrm{EV}_{\mathrm{I}})\) Every set \(\mathbf{V}\in \mathfrak{B}\) is balanced and absorbent.

(EV \(_{II}\) ) For every V \(\in\) B and \(\lambda \neq 0\) in K, we have \(\lambda\) V \(\in\) B (invariance of B under homotheties of non zero ratio).

\((\mathrm{EV}_{\mathrm{III}})\) For every \(\mathbf{V}\in \mathfrak{B}\) , there exists \(\mathbf{W}\in \mathfrak{B}\) such that \(\mathbf{W} + \mathbf{W}\subset \mathbf{V}\) .

Conversely, let E be a vector space on K, and let V be a filter base on E satisfying the conditions \((\mathrm{EV}_{\mathrm{I}})\) , \((\mathrm{EV}_{\mathrm{II}})\) and \((\mathrm{EV}_{\mathrm{III}})\) . Then there exists a topology (and it is unique) on E, compatible with the vector space structure of E, and for which V is a fundamental system of neighbourhoods of 0.

By axiom \((\mathrm{EVT}_{\mathrm{III}}^{\prime})\) we show firstly that the balanced core, \(V_{1}\) , of V, a neighbourhood of 0, is itself a neighbourhood of 0. For there exist \(\alpha > 0\) and a neighbourhood W of 0 such that, if \(|\lambda| \leqslant \alpha\) and \(x \in W\) , then \(\lambda x \in V\) . Since K is non-discrete, there exists \(\mu \neq 0\) in K with \(|\mu| \leqslant \alpha\) and \(\mu W\) is a neighbourhood of 0 for which \(\mu W \subset V\) . Also if \(v \in K\) and \(|v| \leqslant 1\) then \(|v\mu| \leqslant \alpha\) and thus \(v\mu W \supset V\) . Hence \(\mu W \supset V_{1}\) and \(V_{1}\) is a neighbourhood of 0. Also as V is closed so also \(V_{1}\) is closed. Thus the set B of closed balanced neighbourhoods of 0 form a fundamental system of neighbourhoods of 0 in E. By \((\mathrm{EVT}_{\mathrm{I}}^{\prime})\) every neighbourhood of 0 is absorbent; furthermore B satisfies \((\mathrm{EV}_{\mathrm{II}})\) (cf.~I, p.~3, cor.); finally, because of the continuity of \((x, y) \mapsto x + y\) at the point \((0, 0)\) , every fundamental system of neighbourhoods of 0 in E satisfies \((\mathrm{EV}_{\mathrm{III}})\) . The set B satisfies the conditions of the proposition.

Now let E be a vector space over K, and V be a filter base on E satisfying \((\mathrm{EV}_{\mathrm{I}})\) , \((\mathrm{EV}_{\mathrm{II}})\) and \((\mathrm{EV}_{\mathrm{III}})\) . The axiom \((\mathrm{EV}_{\mathrm{I}})\) shows firstly that for all \(V \in V\) , we have -V = V and \(0 \in V\) ; these relations and the axiom \((\mathrm{EV}_{\mathrm{III}})\) show that V is a fundamental system of neighbourhoods of 0, for a topology on E compatible with the additive group structure of E (GT, III, § 1.2). On the other hand the axioms (EVT \(_{I}\) ), (EVT \(_{II}\) ) and (EVT \(_{III}\) ) are immediate consequences of (EV \(_{I}\) ) and (EV \(_{II}\) ), thus the topology defined above satisfies the axiom (EVT), and the proposition is proved.

Remarks. --- 1) In a normed space on a non-discrete valued division ring the set of open balls (resp. closed balls) with centre 0 is a fundamental system of neighbourhoods of 0 which satisfy the conditions \((\mathrm{EV}_{\mathrm{I}})\) , \((\mathrm{EV}_{\mathrm{II}})\) and \((\mathrm{EV}_{\mathrm{III}})\) .

\begin{enumerate}
\def\labelenumi{\arabic{enumi})}
\setcounter{enumi}{1}
\tightlist
\item
  When the division ring of scalars \(\mathbf{K}\) is the field \(\mathbf{R}\) or the field \(\mathbf{C}\) , every filter base \(\mathfrak{B}\) on \(\mathbf{E}\) which satisfies just the two axioms \((\mathrm{EV}_{\mathbf{I}})\) and \((\mathrm{EV}_{\mathrm{III}})\) is a fundamental system of neighbourhoods of 0 for a topology compatible with the vector space structure of \(\mathbf{E}\) . In fact, we need only prove that, in these conditions, for every \(\lambda \neq 0\) in \(\mathbf{K}\) and every \(\mathrm{V} \in \mathfrak{B}\) there exists \(\mathrm{W} \in \mathfrak{B}\) such that \(\lambda \mathrm{W} \subset \mathrm{V}\) . Now from \((\mathrm{EV}_{\mathrm{III}})\) there exists \(\mathrm{W}_{1} \in \mathfrak{B}\) with \(2\mathrm{W}_{1} \subset \mathrm{V}\) , and we deduce, inductively, that for every positive integer \(n\) , there exists \(\mathrm{W}_{n} \in \mathfrak{B}\) such that \(2^{n}\mathrm{W}_{n} \subset \mathrm{V}\) . As \(\mathrm{V}\) is balanced, if we take \(n\) so large that \(2^{n} = |2^{n}| > |\lambda|\) , then \(\mathrm{W} = \mathrm{W}_{n}\) satisfies the condition, as required.
\end{enumerate}

This result does not hold for every non-discrete valued division ring K, for in such a division ring it is no longer necessarily true that \(|me| = m\) for every positive integer m ( \(e\) indicates the unit element of the division ring; cf.~I, p.~22, exerc. 1).

\begin{enumerate}
\def\labelenumi{\arabic{enumi})}
\setcounter{enumi}{2}
\tightlist
\item
  If K is a discrete division ring, conditions \((\mathrm{EVT}_{\mathrm{I}}^{\prime})\) and \((\mathrm{EVT}_{\mathrm{III}}^{\prime})\) are true for any topology on E. Arguing as in prop. 4, one easily sees that if E is a topological vector space on K, then there exists B, a fundamental system of closed neighbourhoods of 0 in E satisfying conditions \((\mathrm{EV}_{\mathrm{II}})\) and \((\mathrm{EV}_{\mathrm{III}})\) . Conversely, if a filter base B on a vector space E over K is such that 0 belongs to all the sets of B and \((\mathrm{EV}_{\mathrm{II}})\) , \((\mathrm{EV}_{\mathrm{III}})\) are true, then B is a fundamental system of neighbourhoods of 0 in a topology compatible with the vector space structure of E.
\end{enumerate}

